\documentclass[11pt]{article}

\usepackage[a4paper,margin=1in]{geometry}
\usepackage{amsmath,amssymb,mathtools}
\usepackage{microtype}

\newcommand{\set}[1]{\left\{#1\right\}}
\newcommand{\card}[1]{\left\lvert#1\right\rvert}
\newcommand{\union}{\mathbin{\cup}}
\newcommand{\logdet}{\operatorname{logdet}}

\title{Federated Semi-Supervised $k$-DPP Selection}
\author{}
\date{}

\begin{document}
\maketitle

\section{Introduction}

A determinantal point process, or DPP, is a probabilistic model for selecting a
subset of items while favoring both useful items and diversity among the
selected items. A k-DPP is a fixed-size version of a DPP and is useful when the
number of items to select is known in advance. In a federated setting, multiple
clients keep their local information private, perform computations locally,
and contribute results that are combined into a shared output. In this project,
each client uses its private context to produce a local k-DPP selection, and
these local selections are combined into a public global set.

\section{Initial Problem Statement}

Let $\mathcal{S}=\set{s_1,\ldots,s_N}$ be a public global universe of
items, and let $\mathcal{C}=\set{1,\ldots,C}$ be a set of clients. Client
$i\in\mathcal{C}$ privately holds an immutable context set
\[
  G_{0,i} \subseteq \mathcal{S}.
\]
The membership, size, and overlap structure of each $G_{0,i}$ are private
client information.

Let $L \in \mathbb{R}^{N \times N}$ be a positive-semidefinite L-ensemble
kernel over $\mathcal{S}$. For any $Y \subseteq \mathcal{S}$,
let $L_Y$ denote the principal submatrix indexed by $Y$.

Each client uses its private context to find a client-specific set $G_{1,i}$.
The initial local objective is
\begin{equation}
  G_{1,i}^{*}
  \in
  \operatorname*{arg\,max}_{\substack{
    G\subseteq\mathcal{S}\setminus G_{0,i}\\
    \card{G}=k}}
  \det\!\left(L_{G_{0,i}\union G}\right).
  \label{eq:local-map}
\end{equation}
Thus, each client selects $k$ items that maximize the joint determinant with
its own private set. The private set $G_{0,i}$ is not shared.

After local optimization, an aggregation operator combines the client-specific
selections into a global public set:
\begin{equation}
  G_1
  =
  \Gamma\!\left(G_{1,1}^{*},\ldots,G_{1,C}^{*}\right)
  \subseteq \mathcal{S}.
  \label{eq:global-aggregation}
\end{equation}

The opertor $\Gamma$ is a public aggregation operator that combines the client-specific selections into a global public set. The aggregation operator is not
necessarily a simple union, but it must be a function of the client-specific selections. The aggregation operator is public and known to all clients.

\section{Baseline}

\subsection{Scope, data, and kernel}

We first study a context-free, two-stage distributed DPP baseline and compare
it with direct centralized sampling. This experiment isolates the effect of
local candidate filtering before global sampling. It does not yet implement
the private semi-supervised problem: all $G_{0,i}=\varnothing$. A client's data partition below is a
candidate pool, not its private context set.

Sample $N=10{,}000$ distinct DNA sequences of length $\ell=20$ uniformly
without replacement from $\{\mathrm{A},\mathrm{C},\mathrm{G},\mathrm{T}\}^{\ell}$.
Reverse complements are treated as distinct
sequences. These items form the public collection $\mathcal{S}$. Randomly
partition it into $C=10$ disjoint, equally sized client pools
$\mathcal{S}_1,\ldots,\mathcal{S}_C$, each containing $1{,}000$ items.
All clients participate; failures and non-IID partitions are not modeled.
$b$ is the final selection size.

Both settings use the same Hamming RBF L-ensemble kernel,
\begin{equation}
  d_H(x,y)=\sum_{r=1}^{\ell}\mathbf{1}\{x_r\ne y_r\},
  \qquad
  L_{xy}=\exp\!\left(-\gamma d_H(x,y)\right),
  \qquad \gamma=0.6.
  \label{eq:baseline-kernel}
\end{equation}
Here Hamming distance is the raw mismatch count, not the fraction of
mismatching positions. The kernel has unit diagonal, no additional quality
weights, and zero jitter. 

\subsection{Distributed baseline: local DPPs followed by a server $k$-DPP}

Each client independently draws one subset $A_i\subseteq\mathcal{S}_i$
from an ordinary, variable-size DPP:
\begin{equation}
  \Pr(A_i=A)
  =\frac{\det(L_A)}{\det(I_{\card{\mathcal{S}_i}}+L_{\mathcal{S}_i})},
  \qquad A\subseteq\mathcal{S}_i.
  \label{eq:baseline-local-dpp}
\end{equation}
This is the L-ensemble law in Kulesza and Taskar~\cite{kulesza2011},
Section~2, Equation~(4). There is no fixed local selection size: both
$\card{A_i}$ and the merged pool size are random. The server receives the
selected global item IDs and forms
\begin{equation}
  P=\bigcup_{i=1}^{C} A_i,
  \qquad m=\card{P}=\sum_{i=1}^{C}\card{A_i}.
  \label{eq:baseline-pool}
\end{equation}
The equality holds because the client pools are disjoint.

For a chosen selection size $b\le m$, the server samples exactly $b$ items from $P$:
\begin{equation}
  \Pr(G_1^{\mathrm{fed}}=Y\mid P)
  =\frac{\det(L_Y)}{e_b(\lambda(L_P))},
  \qquad Y\subseteq P,\quad \card{Y}=b,
  \label{eq:baseline-server-kdpp}
\end{equation}
where $\lambda(L_P)$ denotes the eigenvalues of the pool kernel and $e_b$
is the elementary symmetric polynomial of degree $b$. This is a fixed-size
DPP on the realized pool, following Section~3, Equation~(12)
of~\cite{kulesza2011}. The server returns one shared public set
$G_1^{\mathrm{fed}}$ to all clients. Local selections are only intermediate
candidate sets; their union is not the final output.

Items excluded from $P$ cannot appear in the final server draw.
Consequently, the two-stage law is
not claimed to equal a $k$-DPP on the full collection $\mathcal{S}$.

This is a non-private simulation: the universe and item features are public,
and the server observes client-labeled local selections and final outputs
in plaintext. Communication consists of one upload round and one
broadcast round for the final selection. Other approaches such as SLS, MCMC can be tested against this baseline. 

\subsection{Centralized reference: direct sampling from the full collection}

The centralized setting uses the exact same $10{,}000$ sequences as the
distributed setting, with the same item IDs and kernel
parameters. A single sampler has access to the entire collection and performs
no local selection or pool merging. It directly draws
\begin{equation}
  \Pr(G_1^{\mathrm{cent}}=Y)
  =\frac{\det(L_Y)}{e_b(\lambda(L))},
  \qquad Y\subseteq\mathcal{S},\quad \card{Y}=b.
  \label{eq:baseline-centralized-kdpp}
\end{equation}
Thus, every global item remains eligible, unlike the distributed baseline
where eligibility is restricted to $P$. This reference deliberately permits
a dense global kernel; it is not a scalable private protocol. Both settings
use the same final selection size $b$, which must satisfy $b\le m$ for the
distributed baseline to be feasible.

In the existing implementation, local ordinary DPPs use DPPy's exact
spectral/projection sampler with Gram--Schmidt mode. Both fixed-size stages
(the server and centralized reference) use the same exact rejection sampler,
with Hadamard bounds and complementary-minor identities.
\subsection{Comparison metrics}

Compare centralized and distributed final subsets of the same size using
mean nearest-neighbor distance and its 10th percentile. Both are
\emph{within-subset} measures: for a selected item $x\in Y$, find its
nearest \emph{other} item in the same selected subset $Y$:
\begin{equation}
  d_{\mathrm{NN}}(x;Y)
  =\min_{y\in Y\setminus\{x\}} d_H(x,y),
  \qquad x\in Y,\quad b\ge2.
  \label{eq:baseline-within-nn}
\end{equation}
Distances are raw Hamming mismatch counts, with self-matches excluded.
Compute these distances separately within each method's subset, not between
the two methods' subsets or between the full global collection and a subset.

\paragraph{Nearest-neighbor distance.}
Report the mean of the $b$ within-subset nearest-neighbor distances:
\begin{equation}
  \overline d_{\mathrm{NN}}(Y)
  =\frac{1}{b}\sum_{x\in Y} d_{\mathrm{NN}}(x;Y).
  \label{eq:baseline-mean-nn}
\end{equation}
A higher mean means selected items are farther from their closest neighbors
on average, indicating greater within-subset diversity through local
separation.

\paragraph{10th percentile of nearest-neighbor distance.}
Report the lower-tail summary
\begin{equation}
  q_{0.10}(Y)=\operatorname{Quantile}_{0.10}
    \!\left((d_{\mathrm{NN}}(x;Y))_{x\in Y}\right).
  \label{eq:baseline-nn-tenth-percentile}
\end{equation}
This is the cutoff around the lowest 10\% of the selected items'
nearest-neighbor distances; ties can make the proportion at or below the
cutoff larger than 10\%. A higher value indicates greater separation among
the more closely spaced selected items, and hence greater lower-tail
within-subset diversity. It does not guarantee a minimum distance for every
pair. For both metrics, higher values indicate greater diversity in this
nearest-neighbor sense when subset size and sequence length are held fixed.


\begin{thebibliography}{1}
\bibitem{kulesza2011}
Alex Kulesza and Ben Taskar.
\emph{$k$-DPPs: Fixed-Size Determinantal Point Processes}.
Proceedings of the 28th International Conference on Machine Learning, 2011.
\end{thebibliography}

\end{document}
